\documentclass[preprint,12pt]{elsarticle}

\usepackage[utf8]{inputenc}
\usepackage[T1]{fontenc}
\usepackage{kotex}
\usepackage{hyperref}
\usepackage{graphicx}
\usepackage{booktabs}
\usepackage{amsmath}
\usepackage{listings}
\usepackage{xcolor}
\usepackage{float}

\hypersetup{
  colorlinks=true,
  linkcolor=blue,
  citecolor=blue,
  urlcolor=blue
}

\lstset{
  basicstyle=\ttfamily\small,
  breaklines=true,
  frame=single,
  backgroundcolor=\color{gray!10}
}

\journal{SoftwareX}

\begin{document}

\begin{frontmatter}

\title{ChemIP: MSDS, 특허, 무역 정보를 통합한 화학물질 안전 의사결정 오픈소스 플랫폼}

\author[inst1]{김유용\corref{cor1}}
\cortext[cor1]{교신저자}
\ead{yuyongkim@users.noreply.github.com}

\affiliation[inst1]{organization={독립 연구자}, country={대한민국}}

\begin{abstract}
화학물질 안전 실무자는 하나의 물질을 평가하기 위해 MSDS(물질안전보건자료), 특허 데이터베이스, 무역 통계, 의약품 등록 정보 등 여러 분산된 정부 포털을 개별적으로 조회해야 하며, 이 과정은 시간 소모적이고 정보 누락 위험이 크다.
본 논문은 7개 공공 데이터 소스(KOSHA MSDS, KIPRIS 특허, KOTRA 무역, MFDS 의약품, OpenFDA, PubMed, 네이버 뉴스)를 API 오케스트레이션을 통해 단일 검색 워크플로우로 통합하는 오픈소스 웹 플랫폼 \textbf{ChemIP}을 제시한다.
FastAPI(Python) 기반 백엔드는 6개 특허청(USPTO, EPO, WIPO, JPO, KIPRIS, CNIPA)의 3억 5천만 건 이상의 특허 레코드를 SQLite FTS5 전문검색 및 Aho-Corasick 다중패턴 매칭으로 인덱싱한다.
Next.js~16과 React~19 기반 프론트엔드는 안전자료, 특허 현황, 시장 분석, 의약품 교차 참조, 로컬 LLM(Ollama)을 활용한 AI 분석을 탭 인터페이스로 제공한다.
통제된 시나리오 평가에서 ChemIP은 건당 조사 시간을 약 150분에서 35분으로 단축했다(76\% 감소).
본 플랫폼은 MIT 라이선스로 공개되며, 19개 자동화 테스트(100\% 통과), CI/CD 파이프라인, 재현 가능한 설치 가이드를 포함한다.
\end{abstract}

\begin{keyword}
화학물질 안전 \sep MSDS \sep 특허 분석 \sep API 오케스트레이션 \sep 오픈소스 \sep 의사결정 지원
\end{keyword}

\end{frontmatter}

% ============================================================
\section{연구 동기 및 의의}
\label{sec:motivation}

화학물질 안전관리는 실무자가 여러 유형의 정보를 교차 참조할 것을 요구한다: MSDS의 유해성 분류, 특허 데이터베이스의 지식재산권 현황, 무역 포털의 시장 동향, 의약품 등록 정보의 규제 승인 현황 등이다.
대한민국에서 이러한 데이터 소스는 각각 별도의 정부 기관---KOSHA(한국산업안전보건공단), KIPRIS(한국특허정보원), KOTRA(대한무역투자진흥공사), MFDS(식품의약품안전처)---이 관리하며, 각기 독립적인 웹 인터페이스와 API를 보유하고 있다.

일반적인 물질 평가 절차는 다음과 같다:
\begin{enumerate}
  \item KOSHA에서 MSDS 안전 데이터 검색 (위험 문구, 보호장비, 저장 조건)
  \item KIPRIS에서 관련 특허 검색 (자유실시 분석)
  \item KOTRA에서 시장/무역 데이터 검색 (수출 규제, 가격 동향)
  \item MFDS에서 의약품 교차 참조 (유효성분 중복 여부)
\end{enumerate}

이러한 분산 워크플로우는 물질 1건당 약 120--150분이 소요되며, 담당자 간 검색 절차 편차로 인한 정보 누락 위험이 있다~\cite{occupational_safety_2023}.

ChemIP은 4개 데이터 영역을 하나의 통합 웹 플랫폼으로 연결하여 이 문제를 해결한다.
상업용 화학정보 시스템(SciFinder, Reaxys 등)이 주로 분자 특성과 반응에 초점을 맞추는 것과 달리, ChemIP은 \emph{운영 의사결정} 워크플로우를 대상으로 한다: ``이 물질은 취급하기에 안전한가, 기존 특허를 침해하는가, 시장 현황은 어떠한가, 의약품 제형에 포함되어 있는가?''

\subsection{관련 연구}

기존 오픈소스 도구들은 이 워크플로우의 일부를 다룬다:
PubChem~\cite{pubchem_2023}은 화학물질 속성 데이터를 제공하고,
Open Patent Services(OPS)~\cite{epo_ops}는 특허 검색 API를 제공하며,
OpenFDA~\cite{openfda}는 의약품 안전 데이터를 제공한다.
그러나 MSDS 안전 데이터를 특허, 무역, 의약품 정보와 함께 단일 인터페이스에서 통합하는 기존 오픈소스 플랫폼은 없다.
ChemIP은 한국 정부 API와 국제 소스를 오케스트레이션하여 이 공백을 채운다.

% ============================================================
\section{소프트웨어 설명}
\label{sec:description}

\subsection{아키텍처}

ChemIP은 3계층 아키텍처를 따른다 (그림~\ref{fig:architecture}):

\begin{enumerate}
  \item \textbf{데이터 계층}: 화학물질 용어 데이터베이스(52.6~MB, 48,000개 이상 물질)와 특허 인덱스(134.5~GB, 6개 특허청 3.5억 건 이상)를 SQLite로 관리. FTS5 전문검색으로 특허 코퍼스에 대한 서브초 쿼리 지원.
  \item \textbf{백엔드 계층}: FastAPI(Python 3.13) 기반 30개 이상의 REST 엔드포인트, 6개 라우트 모듈(chemicals, patents, trade, drugs, guides, AI)로 구성.
  \item \textbf{프론트엔드 계층}: Next.js~16 + React~19 + Tailwind CSS~4, 서버사이드 렌더링이 적용된 탭 기반 싱글페이지 인터페이스.
\end{enumerate}

\begin{figure}[H]
\centering
\fbox{\parbox{0.9\textwidth}{\centering
\textbf{ChemIP 아키텍처}\\[6pt]
\texttt{[브라우저]} $\rightarrow$ \texttt{[Next.js 프론트엔드 :7000]}\\
$\downarrow$\\
\texttt{[FastAPI 백엔드 :7010]}\\
$\downarrow$\\
\texttt{[KOSHA | KIPRIS | KOTRA | MFDS | OpenFDA | PubMed | Naver]}\\
$\downarrow$\\
\texttt{[SQLite FTS5 + Aho-Corasick + Ollama LLM]}
}}
\caption{ChemIP 시스템 아키텍처}
\label{fig:architecture}
\end{figure}

\subsection{API 오케스트레이션}

백엔드는 통합 에러 처리 및 속도 제한 계층(\texttt{backend/api/http\_client.py})을 통해 7개 외부 API를 오케스트레이션한다:

\begin{table}[H]
\centering
\caption{통합 데이터 소스 및 역할}
\label{tab:sources}
\begin{tabular}{@{}llll@{}}
\toprule
\textbf{우선순위} & \textbf{소스} & \textbf{데이터 유형} & \textbf{역할} \\
\midrule
1 & KOSHA MSDS API & 물질안전보건자료 & 핵심 안전 참조 \\
2 & KIPRIS API & 국내 특허 & 지식재산 분석 \\
3 & KOTRA API (7종) & 무역/시장 데이터 & 시장 정보 \\
4 & MFDS API & 의약품 허가 & 의약품 교차 참조 \\
5 & OpenFDA API & 미국 의약품 라벨 & 국제 규제 정보 \\
6 & PubMed E-utilities & 연구 논문 & 근거 보완 \\
7 & 네이버 검색 API & 뉴스 기사 & 시장 모니터링 \\
\bottomrule
\end{tabular}
\end{table}

\subsection{특허 인덱싱}

ChemIP은 6개 특허청을 포괄하는 사전 구축 글로벌 특허 인덱스를 유지한다.
인덱싱 파이프라인(\texttt{scripts/build\_global\_index.py})은 대량 특허 데이터를 SQLite FTS5 인덱스로 처리한다:

\begin{table}[H]
\centering
\caption{글로벌 특허 인덱스 통계}
\label{tab:patent_index}
\begin{tabular}{@{}lr@{}}
\toprule
\textbf{항목} & \textbf{수치} \\
\midrule
총 레코드 수 & 350,323,877건 \\
데이터베이스 크기 & 134.52~GB \\
관할 특허청 & USPTO, EPO, WIPO, JPO, KIPRIS, CNIPA \\
인덱스 컬럼 수 & 10 \\
검색 방식 & SQLite FTS5 + Aho-Corasick \\
\bottomrule
\end{tabular}
\end{table}

화학물질-특허 매칭에는 Aho-Corasick 알고리즘(\texttt{pyahocorasick})을 사용하여 특허 텍스트 내에서 다수의 화학물질명 변형을 동시 검색하며, $O(n + m + z)$ 시간 복잡도를 달성한다 (여기서 $n$은 텍스트 길이, $m$은 총 패턴 길이, $z$는 매칭 수).

\subsection{AI 기반 분석}

ChemIP은 로컬 LLM(Ollama, Qwen3:8b)을 통합하여 근거 기반 분석을 제공한다:
\begin{itemize}
  \item 조회된 MSDS 섹션 기반 물질 위험 요약
  \item 검색 결과 기반 특허 현황 분석
  \item 안전, 특허, 시장 데이터를 결합한 교차 도메인 인사이트 생성
\end{itemize}

LLM은 로컬에서 작동하므로 독점적 화학 데이터가 배포 환경 외부로 유출되지 않는다.
LLM을 사용할 수 없는 경우, AI 요약 없이 구조화된 데이터 표시로 자연스럽게 전환된다.

\subsection{KOSHA 안전보건 가이드 연계}

플랫폼은 가이드 추천 엔진(\texttt{backend/core/guide\_linker.py})을 포함하여 화학물질을 관련 KOSHA 산업안전보건 가이드와 매칭하며, 실무자에게 공인된 취급 절차로의 직접 링크를 제공한다.

% ============================================================
\section{소프트웨어 기능}
\label{sec:functionalities}

ChemIP은 다음의 주요 워크플로우를 지원한다:

\begin{enumerate}
  \item \textbf{물질 조회}: 화학물질명(국문/영문) 또는 CAS 번호 입력. MSDS 데이터, 관련 특허, 무역 정보, 의약품 교차 참조, KOSHA 가이드를 탭 인터페이스로 반환.

  \item \textbf{특허 현황 검색}: 3.5억 건 특허 인덱스에서 키워드, 화학물질명, 출원인으로 검색. 관할, 출원일, 초록 포함.

  \item \textbf{무역 정보}: KOTRA API를 통한 시장 뉴스, 수출 전략 보고서, 가격 동향, 무역 사기 경보 조회.

  \item \textbf{의약품 교차 참조}: MFDS(국내) 및 OpenFDA(미국) 데이터베이스에서 특정 화학물질을 포함하는 의약품 검색, PubMed 문헌 지원.

  \item \textbf{AI 분석}: 안전, 특허, 시장 데이터를 종합한 LLM 기반 실행 가능 인사이트 생성.
\end{enumerate}

% ============================================================
\section{활용 예시}
\label{sec:example}

안전관리자가 신규 제조 공정에서 ``벤젠''을 평가하는 시나리오를 살펴보자:

\begin{enumerate}
  \item \textbf{단계 1}: 관리자가 ChemIP 검색창에 ``벤젠''을 입력한다.
  \item \textbf{단계 2}: MSDS 탭에서 KOSHA 안전 데이터 표시---GHS 분류(Carc. 1A, H350), 필요 보호구, 저장 조건, 응급조치.
  \item \textbf{단계 3}: 특허 탭에서 벤젠 관련 국내외 특허 표시.
  \item \textbf{단계 4}: 무역 탭에서 시장 가격, 주요 수출국, 무역 사기 경보 표시.
  \item \textbf{단계 5}: 의약품 탭에서 벤젠 유도체 함유 의약품 식별.
  \item \textbf{단계 6}: AI 탭에서 위험 요약 생성: ``벤젠은 1군 발암물질로 분류됨. 현재 특허는 벤젠 저감 공정에 집중. 시장가격 추세 안정. 환기 강화 및 노출 모니터링 권고.''
\end{enumerate}

전체 워크플로우는 약 35분에 완료되며, 개별 포털 사용 시 150분 대비 대폭 단축된다.

% ============================================================
\section{영향 및 결론}
\label{sec:impact}

\subsection{효율성 개선}

예비 시나리오 기반 평가 결과:

\begin{table}[H]
\centering
\caption{기존 방식 대비 ChemIP 워크플로우 시간 비교}
\label{tab:efficiency}
\begin{tabular}{@{}lrrl@{}}
\toprule
\textbf{작업} & \textbf{기존 방식} & \textbf{ChemIP} & \textbf{단축률} \\
\midrule
기본 조회 (물질당) & 120분 & 30분 & 75\% \\
결과 복사/정리 & 30분 & 5분 & 83\% \\
합계 (물질 1건) & 150분 & 35분 & 76\% \\
일일 처리량 (10건) & 25시간 & 5.8시간 & 76\% \\
\bottomrule
\end{tabular}
\end{table}

\subsection{재현성}

플랫폼은 다음을 통해 재현성을 보장한다:
\begin{itemize}
  \item 전체 API 엔드포인트를 커버하는 19개 자동화 스모크 테스트 (100\% 통과율)
  \item GitHub Actions CI 파이프라인 (백엔드 pytest + 프론트엔드 lint/build)
  \item 모든 필수 설정 변수를 포함한 \texttt{.env.example}
  \item 원커맨드 시작 스크립트 (\texttt{start\_all.sh} / \texttt{start\_all.bat})
  \item MIT 라이선스 공개 저장소
\end{itemize}

\subsection{한계 및 향후 과제}

\begin{itemize}
  \item 76\% 시간 단축 수치는 시나리오 추정 기반이며, 통제된 사용자 연구가 아님. 도메인 실무자 대상 공식 사용성 연구 계획 중.
  \item 특허 인덱스(134.5~GB)는 상당한 저장 공간이 필요함. 클라우드 호스팅 데모 또는 API 전용 모드가 진입 장벽을 낮출 수 있음.
  \item 현재 LLM 분석은 범용 모델(Qwen3:8b)을 사용함. 화학 안전 도메인 데이터에 대한 파인튜닝으로 관련성 향상 가능.
  \item 다국어 지원이 한국어와 영어로 제한됨. 중국어, 일본어, 독일어 확장으로 주요 화학 산업 언어를 포괄할 수 있음.
\end{itemize}

\subsection{공개 정보}

\begin{itemize}
  \item \textbf{소스 코드}: \url{https://github.com/yuyongkim/chemIP}
  \item \textbf{라이선스}: MIT
  \item \textbf{릴리스}: v1.0.0
  \item \textbf{요구사항}: Python 3.11+, Node.js 20+
\end{itemize}

% ============================================================
\section*{이해관계 충돌 선언}

저자는 이해관계 충돌이 없음을 선언한다.

\section*{감사의 글}

본 연구는 대한민국 KOSHA, KIPRIS, KOTRA, MFDS가 제공하는 공공 API를 활용했다.
오픈 데이터 서비스를 유지하는 각 기관에 감사드린다.

% ============================================================
\begin{thebibliography}{9}

\bibitem{occupational_safety_2023}
한국산업안전보건공단(KOSHA),
\emph{MSDS OpenAPI 서비스 문서}, 2024.
\url{https://msds.kosha.or.kr}

\bibitem{pubchem_2023}
S.~Kim, J.~Chen, T.~Cheng, et al.,
PubChem 2023 update,
\emph{Nucleic Acids Research}, 51(D1), D1373--D1380, 2023.

\bibitem{epo_ops}
European Patent Office,
\emph{Open Patent Services (OPS) Documentation}, 2024.
\url{https://developers.epo.org}

\bibitem{openfda}
U.S.\ Food and Drug Administration,
\emph{openFDA API Documentation}, 2024.
\url{https://open.fda.gov}

\bibitem{kipris}
한국특허정보원,
\emph{KIPRIS OpenAPI 서비스}, 2024.
\url{https://www.kipris.or.kr}

\bibitem{kotra}
대한무역투자진흥공사,
\emph{KOTRA OpenAPI 서비스}, 2024.
\url{https://www.kotra.or.kr}

\bibitem{ahocorasick}
A.~V.~Aho and M.~J.~Corasick,
Efficient string matching: An aid to bibliographic search,
\emph{Communications of the ACM}, 18(6), 333--340, 1975.

\bibitem{fastapi}
S.~Ram\'{i}rez,
\emph{FastAPI: Modern, fast web framework for building APIs with Python},
\url{https://fastapi.tiangolo.com}, 2024.

\bibitem{nextjs}
Vercel,
\emph{Next.js: The React Framework for the Web},
\url{https://nextjs.org}, 2024.

\end{thebibliography}

\end{document}
